\section{Preliminaries}
\subsection{Notations}
We use $\prover$ and $\verifier$ to denote
the prover and the verifier.
Let $\setup$ be a generator of bilinear groups, i.e.,
$(p, \g_1, \g_2,$ $\G_1, \G_2, \G_T, e)$ $\leftarrow \setup(\powparam)$,
where all groups have order $p$, and
%Here $\G_1$, $\G_2$, and $\G_T$ all have prime
%order $p$, with 
$\g_1$ ($\g_2$) the generator of $\G_1$ ($\G_2$).
$e: \G_1 \times \G_2 \rightarrow \G_T$ is the bilinear
map.
Let $\F$ be the scalar field of $\G_1$.
$s \pick \F$ represents sampling $s$ uniformly from $\F$.
We use $n$ and $N$ to denote the size of query and lookup tables.
We assume the existence of $n$-th and $N$-th roots of unity 
so FFT is applicable.
%s.t.  for any $a,b \in \Z_p$: $e({\g_1}^a, {\g_2}^b) = e(\g_1,\g_2)^{ab}$ 
%and $e(\g_1,\g_2)$ is the generator of $\G_T$. 
We follow the notations in Groth16 \cite{Groth16}. 
%we write $\G_1$ and $\G_2$ as additive groups.
%That is: given $a\in \Z_p$, we
%denote ${\g_1}^a$ as $[a]_1$, and similarly are group elements in
%$\G_2$ and $\G_T$ denoted. 
For instance, ${\g_1}^a {\g_1}^b$ is written
as $[a]_1 + [b]_1$.
%, $({\g_2}^a)^b$ as $[ab]_2$,
%and $e({\g_1}^a, {\g_2}^b)$ is
%denoted as $[a]_1 \cdot [b]_2$ or $[ab]_T$. 
We use $[n]$ to denote range $[1,n]$.
A vector
$\myvec{t} \in \F^n$ is written as
$(\myvec{t}_1, \ldots, \myvec{t}_n)$.
$\myvec{u} + \myvec{v}$ denotes 
$(\myvec{u}_1 + \myvec{v}_1, \ldots, \myvec{u}_n + \myvec{v}_n)$, 
and $\alpha \myvec{u}$ is $(\alpha \myvec{u}_1, \ldots, \alpha \myvec{u}_n)$.
We use $\myvec{u} \concat \myvec{v}$ to represent
concatenation of $\myvec{u}$ and $\myvec{u}$.
We use $\omega_n$ for the $n$'th root of unity, i.e.,
$(\omega_n)^n = 1$.
We define $\H_n = \{\omega^1_n, \ldots, \omega^n_n\}$,
and the corresponding set of Lagrange base polynomials is denoted as
$\{L_{n,i}(X)\}_{i=1}^n$, where
$L_{n,i}(\omega_{n}^i)=1$ and $L_{n,i}(\omega_{n}^j)=0$ for $j\neq i$.
Given a multi-set $S$, its vanishing polynomial
$z_S(X)$ is defined as $z_S(X) = \prod_{s\in S} (X-s)$.
It is known that $z_{\H_n}(X) = X^n-1$.
Given $s\in \F$, its Pedersen commitment is 
$[s]_1 + r[h]_1$ for a random $r$, where $h$ is unknown to the prover.


Given $\myvec{t} \in \F^n$, its encoding polynomial $t(X)$ has
degree $n-1$ and it is defined as: 
$t(X) = \sum_{i=1}^n \myvec{t}_i L_{n,i}(X)$. 
We use the univariate zk-VPD scheme \cite{zkvsql} to provide
zk-vector commitment. 
Sample $r_t \in \F$, and the masked polynomial
for $\myvec{t}$ is defined as:
$\hat{t}(X) = t(X) + r_t z_{\H_n}(X)$,
and the vector commitment to $\myvec{t}$ is:
$\Commit_{\myvec{t}} = [\hat{t}(s)]_1$, where
$s$ is the trapdoor of a KZG commitment key \cite{Kate10}.
Given a public point $z$, using the univariate port of
the zk-VPD scheme, the prover can convince the verifier
that a Pedersen commitment $\Commit_y$ hides a 
value $y = t(z)$ without disclosing $y$.
Note that $\hat{t}(X)$ is also a 1-leaky masked polynomial
of $t(X)$ \cite[Section 3.6]{CCG23} (also in \cite{GWC19,Lunar}), in
the sense that running standard KZG opening proof for one evaluation
of $\hat{t}(X)$ still keeps $\myvec{t}$ zero knowledge. 2-leaky masking
is similarly achieved by defining
$\hat{t}(X)$ as $\hat{t}(X) = t(X) + (r_{t_1} X + r_{t_0})z_{\H_n}(X)$.
In summary, for any vector $\myvec{u}\in \F^n$, we use $u(X)$,
$\hat{u}(X)$, and $\Commit_{\myvec{u}}$ to denote its
encoding polynomial, masked polynomial, and vector commitment.

We use the notation from
%Camenisch and Stadler in
\cite{Cam97} to specify zk-protocols. 
Consider Schnorr's DLOG as an example:
$\sigdlog(\h)\{(x): \h = [x]_1\}$.
Here $\ttt{DLOG}$ in $\sigdlog$ is a mnemonic.
$(\h)$ is the public information.
The tuple before ``:"
is the secret known by prover only, i.e.,
$(x)$. Then the statement inside curly braces states
the relation: the prover knows the secret discrete logarithm
of $\h$.  
%Applying Fiat-Shamir heuristics  to $\sigdlog(\h)$,
%we get a non-interactive proof and denote it as
%$\prfdlog(\h)$. 
%Schnorr proof can be generalized 
%to multiple bases. 
%We use $1/0 \leftarrow \CheckDLP(\Commit_x, \prf_x, (\g_1,\ldots,\g_k))$
%to denote the verifier function for $\prfdlog$:
% i.e.,
% it returns true if $\prf_x$ is valid for:
%the prover knows $x_1,x_2,\ldots,x_k$ so that
%$\Commit_x = x_1\g_1 + x_2 \g_2 + ... + x_k \g_k$.

% When the context clear,
% $\nu$ and $\omega$ are $n$'th and $N$'th root of unity, 
% i.e., $\nu^n=1$ and $\omega^N=1$. 
% We define $\V = \{\omega^1, \ldots, \omega^N\}$
% and $\H = \{\nu^1, \ldots, \nu^n\}$.
% Let $\{L_i(X)\}_{i=1}^N$ be the Lagrange base polynomials 
% s.t. for each $i\in [n]:$ 
% $L_i(\omega^i)=1$ and $L_i(\omega^j)=0$ for $j\neq i$.
% Similarly, the Lagrange polynomials for $n$ is denoted as
% $l_i(X)$. 
% 
% Similarly, given $\myvec{t}$, 
% $t(X) = \sum_{i=1}^n \myvec{t}_i \cdot l_i(X)$.

\subsection{Commitments}
We use a variety of commitments in this paper. We first provide
a formal definition.
\begin{definition} Let $\myvec{v} \in \F^n$ be a vector of field elements.
A commitment scheme $\com$ is a set of algorithms
$(\comsetup, \comprove, \comverify)$, where:
\begin{enumerate}
 \item $(\pk,\vk) \leftarrow \comsetup(1^\lambda)$. On input security parameter
 $\lambda$ generates prover and verifier keys $(\pk, \vk)$.
 \item $\com \leftarrow \comcommit(\pk, \myvec{v} \in \F^n, r)$. Generates
 a succinct commitment to vectof $\myvec{v}$, where $r$ is a hiding factor.
 \item $\prf \leftarrow \comprove(\pk, \myvec{v}, r, y)$. Generates
 a proof $\prf$ which shows that (polynomial built from) $\myvec{v}$ 
 evaluates to $y$ at $r$, or $\myvec{v}[r] = y$ depeding on the context. 
 \item $1/0 \leftarrow \comverify(\vk, \com, \prf, r, y)$. Verify
 the above proof.
\end{enumerate}
\end{definition}
We require $\com$ satisfies properties such as {\em complete hiding}: 
one cannot infer any information about the vector elements from a commitment,
{\em computational binding}: it is computationally hard to find 
$\myvec{v}' \neq \myvec{v}$ s.t. 
$\comcommit(\myvec{v}) = \comcommit(\myvec{v}')$.

In this paper, we consider three forms of vector commitments:
Pedersen vector commitment ($\pedcom$), KZG commitment
($\kzgcom$), and Lagrange basis vector commitment ($\veccom$). 
Given a vector $\myvec{a}$, we write its KZG commitment as
$\barkzg{\myvec{a}}$. Similarly its Pedersen vector commitment
is denoted as $\barped{\myvec{a}}$ and its
Lagrange basis vector commitment $\barvec{\myvec{a}}$.
We note that using the Quasi-linear 
NIZK (QA-NIZK) technique \cite{Gonzalez15}, as shown in
Lego-Snark \cite{LegoSnark}, one can prove the equivanece
of these commitments encoding the same vector.
The default commitment scheme (as used in Nova/SuperNova)
is the Pedersen Vector Commitment. By default $\bar{\myvec{a}}$
denotes $\barped{\myvec{a}}$, if otherwise specified.

\subsubsection{Adding Zero Knowledge to KZG}
\noindent{\bf Bounded Leaky-Zk}
We first discuss the ``leaking-zk".
The idea is to mask a polynomial with a random poly.
It is formally defined as bounded-zk in Lunar \cite{Lunar},
but was unofficially used in PlonK \cite{GWC19}\footnote{ 
NOTE: It was mentioned in Lunar's
introduction in bounded-leaky zk, 
however, didn't find where it is in \cite{GWC19}. {\bf CHECK LATER.}}
and then summarized in \cite{CCG23} (section adding zero
knowledge). We formally defined it below.

\begin{lemma} Let $\myvec{f} \in \F^n$, and let 
$f(X) = \sum_{i=1}^n \myvec{f}_i \cdot l_i(X)$.
Given a fixed $b>0$, we say that $\hat{f}$
is a {\em $b$-bounded-zk extension} of $f(X)$ 
in the sense that there exists a $b$-1 degree
random polynomial $R(X)$ s.t.
\[
	\hat{f}(X) = f(X) + R(X) z_{\H}(X)
\]
$\hat{f}$ satisfies the following properties:
\begin{enumerate}
	\item For $i\in [n]:$ $\hat{f}(\omega^i) = \myvec{f}_i$.
	\item Given $b$ evaluations of $\hat{f}$ at random points,
		for any random vector $\myvec{g} \in \F^n$,
		there exists a $b$-1 degree poly $U$ s.t.
		$\hat{f}(X)$ = $g(X)$ + $U(X) z_{\H}(X)$.
\end{enumerate}
\end{lemma}

Property (2) essentially states that as long as there are no more than
$b$ evaluations, $\hat{f}$ can be generated from any arbitrary
vector, thus leaking no information about $\myvec{f}$.
Note that for the $1$-leaky case, $R(X)$ is degree $0$.

\subsubsection{Complete-Zk KZG (NOT USED TO DEL)}
An alternative way is to applying Pedersen commitment style blinding
factor and use Schonorr style zk-protocol to cancel out the blinding
factor. This is done in \cite{binacc22} and our own work
\ttt{zkregex}.  The \ttt{zkregex} construction is an improved  univariate
variation of the ZK-VPD construction in \cite{zkvsql}.

This approach guarantees full zk at a small $O(1)$
cost in proof size and verifier cost. We officially call
it $\protokzg$ which defines the following operations:
\begin{enumerate}
	\item $\srs \leftarrow \setup(N, \secparam):$ the trusted
		set up generates the CRS, which essentially
		are KZG prover key
		$\left( ([s^i]_1)_{i=1}^N, [z_\V(s)]_1, [z_\H(s)]_1 \right)$.
	\item $\Commit_f \leftarrow \commit(f, r, \srs)$: 
		we have two cases: (1) $f\in \F_{<N}[x]$: generate
		compute $\Commit_f = [f(s)]_1 + r_f[z_\V(s)]_1$;
		(2) $f\in \F_{<n}[x]$: compute
		$\Commit_f = [f(s)]_1 + r_f[z_\H(s)]_1$.

	\item $(\Commit_v, \prf) \leftarrow \prove(f, r_f, t, v, r_v, r_\prf, 
		\srs)$:
		generate the valuation proof that $f$ evaluates to some secret
		value $v$ at point $t$ and $\Commit_v$ commits to $v$.
		$r_f$, $r_v$, $r_\prf$ are the blinding openings for the
		Pedersen commitments to $f$, $v$, $\prf$ respectively.
	\item $1/0 \leftarrow \verify(\Commit_f, t, \Commit_v, \prf, \srs)$:
		verify the evaluation proof is valid.
\end{enumerate} 

We recall the details of $\protokzg$ in \ttt{zkregex}.
$\prover$ computes $q(X) = (f(X) - v)/(X-t)$ and
$\Commit_q = [q(s)]_1 + r_q[z_\V(s)]_1$ for random
$r_q$. $\prover$ also computes a balancing term $B$:
\[
	B = (r_f - r_v) e([1]_1, [z_\V(s)]_2)
		 - r_q e([z_\V(s)]_1, [s-t]_2)
\]
the proof is tuple $(\Commit_q, B, \prfdlp(B, (e([1]_1, [z_\V(s)]_2), e([1]_1, [s-t]_2))$ s.t.
\[
	e(\Commit_f - \Commit_v, [1]_2) = 
	e(\Commit_q, [s-t]_2) + B
\]
$\verifier$ performs the pairing check and the $\prfdlp$ check.
Note that because $\Commit_v$ is used, the construction
is complete-zk.

{\bf Remark:}  we specifically note that
$\Commit_{\myvec{T}}$ can be regarded as a standard
KZG commitment to $\hat{T}(X)$,  as well as a
$\protokzg$ commitment to $T(X)$ at the same time.

\section{Basic Crypographic Constructions}
In this section we list some basic constructions that we need.
Technical details can be found in Appendix.

Given a vector $\Commit_{\myvec{v}}$, and a Pedersen
commitment $\Commit_u$, $\prfsum(\Commit_{\myvec{v}}, r, \Commit_u)$
proves that $u = \sum_{i=1}^N v_i r^i$.
$\prfprod(\Commit_\myvec{u}, \Commit_\myvec{v}, \Commit_\myvec{w})$
proves that for each $i$: $\myvec{w}_i = \myvec{u}_i \times \myvec{v_i}$.

$\prfperf(\Commit_{\myvec{u}}, \Commit_{\myvec{v}})$ to prove that
$\myvec{v}$ is a permutation of $\myvec{u}$. This can leverage
the grand-product used in Plonk system (details later).

$\prfset(\Commit_\myvec{u}, \Commit_\myvec{s})$: to prove that
vector $\vec{s}$ encodes the the standard set of elements from
$\myvec{u}$. The basic idea is simple: first produce
a vector $v$ which is a sorted permutation, run $\prfperf$ 
first, and then prove $\myvec{v}$ is sorted, and 
then run $\prfpack$ to produce $\myvec{s}$.



\subsection{Lookup Argument}
% We need homomorphic and zero knowledge look-up arguments \cite{plookup,Caulk22,CaulkPlus,
% Baloo,Flookup,cq} in the sense that: (1) it is zero knowledge, and (2)
% it provides an additional Pedersen commitment to the array of locations
% of query elements in lookup tables.
% We call it zero knowledge lookup argument enhanced with
% aux information about positions (zk-aux-lookup),
% and denote it as $\zkauxlookup$. 
% The construction can be adapted from
% the segment lookup argument in \cite{CCG23}.
% %For the ``standard" zk-version
% %without the aux information on locations, we denote it as
% %$\zklookup$. 
% 
% \begin{enumerate}
% 	\item $\srs \leftarrow \lsetup(N, \secparam): $ a trusted set-up
% 	given security parameter $\secparam$ and container table size
% 	limit $N$, samples bilinear groups and generates common reference
% 	string $\srs$. We here do not distinguish between prover and
% 	verifier keys, however, in general verifier key is much smaller.
% 
% 	\item $(\aux,\Commit_{\myvec{T}}) 
% 		\leftarrow \preprocess(\myvec{T})$: Given the container
% 	table $\myvec{T}$, it generates preprocessed information $\aux$ and
% 	$\Commit_{\myvec{T}}$, a vector commitment to $\myvec{T}$.
% 
% 	\item $(\Commit_{\myvec{l}}, \prf) 
% 	\leftarrow \prove(\crs, \aux, \myvec{T}, \myvec{t})$: produces
% 	a zero knowledge proof $\prf$ which asserts that each element in
% 	$\myvec{t}$ belongs to $\myvec{T}$. The aux information $\aux$ is
% 	used for speeding up proof generation (e.g., in $\cq$ \cite{cq}).
% 	Let $\myvec{l}$ descries the location of each 
% 	$\myvec{t}_i$ in $\myvec{T}$, i.e., $\forall i\in [1,n]:$
% 	$\myvec{l}_i = \myvec{T}_j$ s.t. $\myvec{T}_j=\myvec{t}_i$.
% 	$\Commit_{\myvec{l}}$ is a Pedersen commitment to 
% 	$\myvec{l}_i$.	
% 	\item $1/0 \leftarrow \verify(\crs, \Commit_{\myvec{T}}, 
% 	\Commit_{\myvec{t}}, \Commit_{\myvec{l}}, \prf):$
% 	checks that the zk-aux-lookup argument is valid.
% \end{enumerate}
% 
% The definition for completeness, knowledge soundness and zero knowledge
% is standard and will be presented in a full-version of this manuscript.
% The difference from a zero knowledge subset proof is that
% $\myvec{t}$ is a multi-set (allowing repititions).
% This is needed later for our aggregated non-membership proof.
% The construction can be adapted from the segment lookup
% argument from Sublonk \cite{CCG23}. The observation is that
% it is possible to embed location information by extending
% the log-derivative lemma introduced in 
% \cite{Hab22}.
% 
% \begin{lemma} 
% \label{lemma:ext}
% (Segment-Lookup Lemma in \cite{CCG23}, Extension of Lemma 5 \cite{Hab22}  and Lemma 2.4 \cite{cq}) 
% Given $f \in \F^n$ and $t \in \F^N$, and $\omega^N = 1$.
% We have $f \subseteq t$ as sets if and only if 
% for some $m \in \F^N$ and $p \in \F^n$, for two randomly
% sampled $X,Y \in \F$, the following holds:
% \[
% 	\Pr \left[
% 		\sum_{i\in [n]} \frac{1}{X + f_i + Y \cdot p_i} \neq
% 		\sum_{i\in [N]} \frac{m_i}{X + t_i + Y \cdot \omega^i}  
% 	\right] = \negli
% \]
% \end{lemma}
% 
% Here $p(i)$ captures the position of $\myvec{t}_i$ 
% in $\myvec{T}$. Then, by enforcing the new equation in an adapted
% $\cq$ protocol we are able to embed location information.
% The protocol also adds zero-knowledge. The details are presented
% in Appendix \ref{apdx:cq}.
% 

Let $\myvec{t} \in \F^n$ and
$\myvec{T} \in \F^N$, a lookup argument for $\myvec{t} \lookup \myvec{T}$
asserts that for each $i\in [n]$: $\myvec{t}_i \in \myvec{T}$.
Note that both $\myvec{t}$ and $\myvec{T}$ may contain duplicates.
We need a homomorphic and zero knowledge look-up argument, and
a slight zk-enhancement of $\cq$ \cite{cq} satisfies our needs.
The basic idea of $\cq$ is to pre-compute commitments of
quotient polynomials in $O(N\log(N))$ time. Then 
for arbitrary query table $\myvec{t}$
of size $n$, the prover work is $O(n\log(n))$, thus 
independent of $N$. We need to enhance $\cq$ so that it is zero knowledge.
In our extension, we exploit bounded leaky-zk as in \cite{Marlin,CCG23},
and Schnorr style $\Sigma$-protocols as in \cite{binacc22}. 
We will present its full technical details in the extended version 
of this paper.  Similar constructions for the zk-enhancement
of $\cq$ can be found in two concurrent work:
segment lookup \cite{CCG23} and matrix lookup \cite{MatrixLookup}.
%For the ``standard" zk-version
%without the aux information on locations, we denote it as
%$\zklookup$. 
We denote it as $\zkauxlookup$, as shown below:

\begin{enumerate}
	\item $(\pk,\vk) \leftarrow \lsetup(N, \secparam): $ a trusted set-up
	given security parameter $\secparam$ and lookup table size
	limit $N$, samples bilinear groups and generates the
	prover and verifier keys $(\pk, \vk)$ for KZG.

	\item $(\aux_{\myvec{T}},\Commit_{\myvec{T}}) 
		\leftarrow \preprocess(\myvec{T}, \pk)$: Given the lookup 
	table $\myvec{T}$, it generates 
	$\aux_{\myvec{T}}$ (the preprocessed information)
	and $\Commit_{\myvec{T}}$ (a commitment to $\myvec{T}$).

	\item $\prf 
	\leftarrow \prove(\pk, \aux_{\myvec{T}}, \myvec{T}, \myvec{t})$: produces
	a zero knowledge proof $\prf$ for
	$\myvec{t} \lookup \myvec{T}$.

	\item $1/0 \leftarrow \verify(\vk, \Commit_{\myvec{T}}, 
	\Commit_{\myvec{t}}, \prf):$ checks that $\prf$ is valid.

	\item $\prf 
	\leftarrow \foldprove(\pk, \aux_{\myvec{U}}, \myvec{U}, 
		\aux_{\myvec{V}}, \myvec{V}, 
		\myvec{u}, \myvec{v}, \alpha)$: produces
	a zero knowledge proof $\prf$ for
	$\myvec{u} + \alpha \myvec{v} \lookup 
		\myvec{U} + \alpha \myvec{V}$.

\end{enumerate}

% The definition for completeness, knowledge soundness and zero knowledge
% is standard and will be presented in a full-version of this manuscript.
% The difference from a zero knowledge subset proof is that
% $\myvec{t}$ is a multi-set (allowing repetitions).
% Also the $\foldprove$ operation allows to easily prove
% multi-column tables exploiting the homomorphic property
% of $\cq$. There are two options to provide zk. One way is
% bounded leaky-zk \cite{Marlin,CCG23} where the degree of a polynomial
% is increased by $k$ if it is known that it will be evaluated
% at most $k-1$ times. Another alternative is to blind each
% KZG commitment with a random factor, which increases
% the proof with $O(1)$ size extra Schnorr style proof. 
% In the rest of the paper, for convenience of presentation,
% all protocols can be enhanced with zk, but we omit the details.

\begin{theorem} Under the $Q$-DLOG assumption \cite{cq} in
the Algebraic Group Model (AGM) and Random Oracle Model (ROM),
$\zkauxlookup$ is perfectly complete, computational
knowledge sound, and zero knowledge.
\end{theorem}

The proof generally follows the analysis of 
$\cq$ \cite{cq}, and is similar to that of 
\cite{CCG23,MatrixLookup}. We delay the details 
to the extended version.

Based upon $\zkauxlookup$,
we develop $\prfrange(\Commit_{\myvec{a}}, 2^B)$,
a batched zk-range proof which asserts 
that $\Commit_{\myvec{a}}$ is
a Pedersen vector commitment to $\myvec{a} \in \F^n$,
and each $\myvec{a}_i$ in 
$\myvec{a}$ is in range $[0, 2^B)$. 
It has $O(1)$ proof size and $O(n\log(n))$
 prover complexity.


\subsection{Commit-and-Prove zkSNARK}
In our application scenario, we have specialized $\Sigma$-protocol
(e.g., lookup arguments), and there are boolean and arithmetic
logic which can be more conveniently captured by R1CS and handled
by general zk-SNARK systems such as Groth16 \cite{Groth16}.
We replicate the Commit-and-Prove scheme in LegoSnark \cite{LegoSnarkFull},
and bind a Groth16 system with the zk-cq lookup argument, 
formally defined as below.

\begin{definition}  A Commit-and-Prove for Lookup SNARK System (CPL-SNARK)
is a collection of algorithms that supports the argument for
$S = (R=(\myvec{io},\myvec{w} \cup \{r\},A,B,C), \vec{f}, 
\{(\myvec{v}_i, T_i)\}_{i=0}^k)$, 
where $R$ is an R1CS instance. 
$\vec{f} \subseteq \vec{w}$ is a subset of variables to be
committed before $r$ of $R$ is determined.
For each $i$, $\myvec{v}_i$ is a subset of $\myvec{w}$.
\begin{enumerate}
 \item $(\pk, \vk) \leftarrow \setupcpl(\lambda, S)$.
 \item $\prf \leftarrow \prvcpl(\pk, S)$. It asserts the
 hidden $\myvec{w} \cup \{r\}$ is a valid witness of
 $R$, and $r = hash(\vec{f})$, and
 for each $\myvec{v}_i$, all of its elements
 is contained in table $T_i$.
 \item $0/1 \leftarrow \vercpl(\vk, io)$.
\end{enumerate}
\end{definition}

We note that the $r$ is used to perform fingerprinting check
of two relations. For example given two vectors $\myvec{a}$
and $\myvec{b}$ we could have an arithmetic circuit, which given
a random challenge $r$ and test $\prod_{i=1}^{n} (\myvec{a}_i-r) 
= \prod_{i=1}^n (\myvec{b}_i-r)$. However, this test is only
valid when $\myvec{a}$ and $\myvec{b}$ is ``fixed" before
random challenge $r$ is given as the input to the circuit.

{\bf TO BE EXPANDED:} Given that there are multiple circuits to
be called depending on the first stage, and the number of circuits
to be deployed ranging between $0$ to $64$, we have several alternative
approaches: (1) aggregation of proofs (2) folding scheme (with
or without IVC). Aggregation has $\log(n)$ proof size
and folding provides $O(1)$ proof size, and much faster prover
performance (as no need to prove each indidivual instance just
fold them with a couple of MSMs).  The problem is that
IVC based folding usually needs to encode folding verifier circuit,
which typically requires a cycle of curves, and for preformance issue
(e.g., Pasta curves) do not provide pairing, which is needed for
the CQ protocol for large tables. In this case, we mix several
ingredients together:
\begin{enumerate}
 \item We apply folding to all instances of circuits.
 \item We use commit-and-prove scheme to extract out the
 	input and output of circuits, as well as the query table
	inside circiut, using QA-NIZK.
 \item We use GIPA to aggregate folding.
\end{enumerate}

This eventually yields a $\log(k)$ (given $k$ is the number of
instances) proof size, and only one curve needed, for
folding scheme (even IVC).

